\honest{My Bayesian Enlightenment}{Eliezer Yudkowsky}{2008}

I remember, dimly, the first time I called myself a Bayesian. Someone asked: if a mathematician says ``I have two children, and at least one of them is a boy,'' what is the probability that both are boys? I pointed out that a mother of a boy and a girl could as well have mentioned the girl, and I got 1/2. \nb{The answer is right under that assumption. It is the standard analysis since Martin Gardner granted that his version of the puzzle was ambiguous. In the well-posed version the post also gives, you ask the question, and the answer is 1/3.}

Someone told me that 1/2 was the Bayesian answer, and that in ``orthodox statistics'' the answer is 1/3, because frequentists do not guess what the mathematician would say. I replied that they were assuming that probability to be 1, and asked: ``How can there possibly be such a thing as non-Bayesian statistics?'' That is how I learned that I was a Bayesian, and as far as I can tell I was born that way: what Bayesians said seemed the obvious way to do it, and what frequentists said sounded like ``the elaborate, warped, mad blasphemy of dreaming Cthulhu.'' \nb{The two answers differ in how the statement came to be made, not in the school of the person computing. A frequentist who counts the families in which a parent would say this sentence also gets 1/2, if parents of one of each mention either sex equally often.}

That was not my enlightenment. I marked Bayes off as obvious and went no further than the rule itself. I still thought of probability theory as a tool, not a law; like nearly all AGI hopefuls, I collected techniques and algorithms, looking for tools instead of understanding, and Bayes's Rule was a neat one. Nor was my enlightenment heuristics and biases. I met those results on a web page made from a PowerPoint introduction to behavioral economics, which gave no references. I emailed the author to ask whether the experiments were real, and he sent me a scan of Tversky and Kahneman's 1973 paper. Even then I put it on a list; I thought I could get along without anything not online. Emil Gilliam, annoyed by that theory, bought me \textsc{Judgment Under Uncertainty}. It showed me that my Traditional Rationality was inadequate, that there was much more than doing what Feynman told you, and it held up Bayes as the gold standard. But that was not all the way.

My memory for the order of events in everyday life is poor, though I remember causal structure well.

I am not sure whether I read E.~T. Jaynes's \textsc{Probability Theory: The Logic of Science} before or after the day I saw the magnitude of my folly. But Jaynes did it. There was probability as ``The Rules'', ``inviolable on pain of paradox'', with the same answer by every legitimate route. If you approximated the Rules because they cost too much to compute, you would still do less than optimal, however necessary the compromise. Jaynes would take the different answers others had reached and trace each to the illegitimate step. Not an answer, but the answer. Having looked back on all the vague answers that had led me into paradox, I saw the level above mine, and could no longer picture building an AI on such answers and surviving. \nb{This matches Jaynes's own claims, that the rules do ``uniquely and optimally'' what orthodox statistics does with ``intuitive ad hoc devices.'' How far the rules are unique is discussed in the notes on ``Beautiful Probability''.} The AGI hopefuls' dreams of Friendliness, ``Like frequentist statistical methods'', did not agree with one another; many were invented on the spot when I asked. I had studied the problem for years and knew what their plans would run into. ``I don't see why this would fail'' only reflects your own ignorance, and I would be doomed too if I settled for what seemed like a good idea. What it would take was something like the Jaynes level: not ``here's my bright idea'' but the only correct way to do it, and why. \nb{Frequentist theory derives its methods from stated criteria and proves results about them, such as the Neyman--Pearson lemma of 1933.}

A passage flashed through my mind: ``Do nothing because it is righteous, or praiseworthy, or noble, to do so; do nothing because it seems good to do so; do only that which you must do, and which you cannot do in any other way.'' \nb{The post does not say whose words these are. Goodreads gives them as Ursula K. Le Guin's, from \textsc{The Farthest Shore}.} So I called a full stop. I held my Friendly AI designs to that standard. None of my old theories met it, came close, or was on a track toward it, so I threw them all out and took up probability theory and decision theory, hoping to extend them to reflectivity and self-modification. Seeing cognition as Bayes-structure, my naturalistic awakening, and seeing that Traditional Rationality was not strict enough were other parts of the same enlightenment.

Later, reading Judea Pearl's \textsc{Probabilistic Reasoning in Intelligent Systems}, I found that precision saves time. I had once worked on nonmonotonic logics, and I saw how much time I would have wasted on ad hoc systems without Pearl's key; the savings are measured in careers, not months. Only by holding to a higher standard of precision had I begun to think at all about many important issues. Precision is not formality, or inventing a new logic to throw at a problem. People who call it impossible say so ``because human beings are lazy,'' without trying for five minutes. \nb{The post gives no evidence for this cause.}

Without an inconveniently high standard, like a proof in which one wrong step can carry you anywhere, you will not chase down the tiny notes of discord that lead to new concerns. Finding a standard high enough to make you start thinking is itself hard.

``Oops'' is something to look forward to. It means your present self is a drooling imbecile, but also that your future self will gain powers your present self does not dream of, and that you have not yet passed your peak. The scream of embarrassment at your younger self is ``the sound that rationalists make when they level up.'' Sometimes I worry that I am not leveling up as fast as I used to, and do not know whether I am getting the hang of things or my neurons are slowly dying.
